\input{estilo}
% En la pizarra la flecha va de la lectura de T2 a la escritura de T1: T2 ve a
% T1. Va en verde, como la lectura correcta de lector-atrasado. Las guías grises
% bajan al eje el commit de T1 y el comienzo de T2: el "antes" es en tiempo
% físico, como dice el texto de esa sección, y todavía no hay timestamps.
\begin{document}
\begin{figura}\begin{tikzpicture}
  \tx{1.2}{1}  \tx{0}{2}
  \transaccion[w]{1.2}{1.4}{4.8}{\op{W}{x}}
  \transaccion[r]{0}{5.6}{9.2}{\op{R}{x}}
  \draw[->, segundo] (r.north) .. controls +(0,1.1) and +(1.3,0.7) .. (w.east)
    node[pos=0.2, right=2pt, etiq, text=segundo] {\var{T}\textsubscript{2} ve a \var{T}\textsubscript{1}};
  % el antes, en tiempo físico
  \draw[densely dotted, draw=suave, line width=0.5pt] (4.8,1.2) -- (4.8,-0.8) (5.6,0) -- (5.6,-0.8);
  \hora[below left]{(4.8,-0.8)}{commit \var{T}\textsubscript{1}}
  \hora[below right]{(5.6,-0.8)}{empieza \var{T}\textsubscript{2}}
  \tiempo[tiempo físico]{1.4}{9.8}{-0.8}
\end{tikzpicture}\end{figura}
\end{document}
